% Partie : set-theory
% Chapitre : spine
% Section : zf

\documentclass[../../../include/open-logic-section]{subfiles}

\begin{document}

\olfileid[fr]{sth}{spine}{zf}
\olsection{\texorpdfstring{$\Z$}{Z} et \texorpdfstring{$\ZF$}{ZF} : un jalon}

Avec la fondation, nous franchissons un autre jalon important. Nous
avons étudié les théories $\Zminus$ et $\ZFminus$, dont nous disions
qu'il leur manquait quelque chose. Ce quelque chose est la fondation.
Ainsi :

\begin{defn}
La théorie $\Z$ ajoute la fondation à $\Zminus$. Elle a donc pour
axiomes l'extensionnalité, l'union, la paire, l'ensemble des parties,
l'infini, la fondation et toutes les instances du schéma de
séparation.

La théorie $\ZF$ ajoute la fondation à $\ZFminus$. Autrement dit,
$\ZF$ ajoute à~$\Z$ toutes les instances du schéma de remplacement.
\end{defn}

Une question a pourtant pu vous venir à l'esprit. Si la régularité
est équivalente à la fondation dans $\ZFminus$, et que la
justification de la régularité est claire, pourquoi ce détour ?
Pourquoi retenir la fondation, plutôt que la régularité, comme axiome
de base ?

Si l'on laisse de côté les raisons historiques (qui tiennent à savoir
qui a formulé quoi, et quand), la raison essentielle est que la
fondation peut être énoncée sans faire appel à la définition des
$V_\alpha$. Cette définition reposait sur tout le travail de la
\olref[recursion]{sec} : il a fallu démontrer le théorème de
récursion transfinie pour la justifier. Or notre démonstration de ce
théorème utilisait le \emph{remplacement}. Ainsi, bien que la
fondation et la régularité soient équivalentes modulo $\ZFminus$,
elles ne le sont pas modulo $\Zminus$.

La situation est même plus radicale que ne le laisse penser cette
remarque. Bien que cela dépasse largement le propos de ce livre, il
s'avère que $\Zminus$ et $\Z$ sont toutes deux trop faibles pour
définir les $V_\alpha$. Si nous ne travaillons que dans $\Z$, la
régularité, telle que nous l'avons formulée, n'a donc même pas de
\emph{sens}. C'est pourquoi notre axiome officiel est la fondation,
et non la régularité.

Désormais, sauf indication contraire, nous travaillerons dans $\ZF$
sans autre précision.

\end{document}
